\documentclass[10pt,a4paper]{amsart}
\usepackage{amsmath,amssymb,amsthm,graphicx}
\usepackage{verbatim}
\newenvironment{reference}{\comment}{\endcomment}
%\newenvironment{reference}{}{}
\newenvironment{slogan}{\comment}{\endcomment}
\newenvironment{history}{\comment}{\endcomment}

% For commutative diagrams we use Xy-pic
\usepackage[all]{xy}

% We use 2cell for 2-commutative diagrams.
\xyoption{2cell}
% Theorem environments.
%
\theoremstyle{plain}
\newtheorem{theorem}[subsection]{Theorem}
\newtheorem{proposition}[subsection]{Proposition}
\newtheorem{lemma}[subsection]{Lemma}

\theoremstyle{definition}
\newtheorem{definition}[subsection]{Definition}
\newtheorem{example}[subsection]{Example}
\newtheorem{exercise}[subsection]{Exercise}
\newtheorem{situation}[subsection]{Situation}

\theoremstyle{remark}
\newtheorem{remark}[subsection]{Remark}
\newtheorem{remarks}[subsection]{Remarks}

\numberwithin{equation}{subsection}

% Macros
%
\def\lim{\mathop{\mathrm{lim}}\nolimits}
\def\colim{\mathop{\mathrm{colim}}\nolimits}
\def\Spec{\mathop{\mathrm{Spec}}}
\def\Hom{\mathop{\mathrm{Hom}}\nolimits}
\def\Ext{\mathop{\mathrm{Ext}}\nolimits}
\def\SheafHom{\mathop{\mathcal{H}\!\mathit{om}}\nolimits}
\def\SheafExt{\mathop{\mathcal{E}\!\mathit{xt}}\nolimits}
\def\Sch{\mathit{Sch}}
\def\Mor{\mathop{\mathrm{Mor}}\nolimits}
\def\Ob{\mathop{\mathrm{Ob}}\nolimits}
\def\Sh{\mathop{\mathit{Sh}}\nolimits}
\def\NL{\mathop{N\!L}\nolimits}
\def\CH{\mathop{\mathrm{CH}}\nolimits}
\def\proetale{{pro\text{-}\acute{e}tale}}
\def\etale{{\acute{e}tale}}
\def\QCoh{\mathit{QCoh}}
\def\Ker{\mathop{\mathrm{Ker}}}
\def\Im{\mathop{\mathrm{Im}}}
\def\Coker{\mathop{\mathrm{Coker}}}
\def\Coim{\mathop{\mathrm{Coim}}}

% Boxtimes
%
\DeclareMathSymbol{\boxtimes}{\mathbin}{AMSa}{"02}

%
% Macros for moduli stacks/spaces
%
\def\QCohstack{\mathcal{QC}\!\mathit{oh}}
\def\Cohstack{\mathcal{C}\!\mathit{oh}}
\def\Spacesstack{\mathcal{S}\!\mathit{paces}}
\def\Quotfunctor{\mathrm{Quot}}
\def\Hilbfunctor{\mathrm{Hilb}}
\def\Curvesstack{\mathcal{C}\!\mathit{urves}}
\def\Polarizedstack{\mathcal{P}\!\mathit{olarized}}
\def\Complexesstack{\mathcal{C}\!\mathit{omplexes}}
% \Pic is the operator that assigns to X its picard group, usage \Pic(X)
% \Picardstack_{X/B} denotes the Picard stack of X over B
% \Picardfunctor_{X/B} denotes the Picard functor of X over B
\def\Pic{\mathop{\mathrm{Pic}}\nolimits}
\def\Picardstack{\mathcal{P}\!\mathit{ic}}
\def\Picardfunctor{\mathrm{Pic}}
\def\Deformationcategory{\mathcal{D}\!\mathit{ef}}
\usepackage{cleveref}
\crefname{originallemma}{original lemma}{original lemmas}
\Crefname{originallemma}{Original lemma}{Original lemmas}
\makeatletter
\newlabel{spaces-properties-lemma-map-into-reduction}{{03JJ}{source}{}{sourcepointer-1}{}}
\newlabel{spaces-properties-lemma-map-into-reduction@cref}{{[originallemma][0][]03JJ}{[1][0][]source}{}{}{}}
\makeatother
\title{A rational map from a reduced algebraic space to a separated one is represented by a morphism on its domain of definition}
\author{The Stacks Project Authors}
\date{stacks-project:spaces-morphisms.tex:local-snapshot-20261012; GNU FDL 1.2 or later}
\begin{document}
\maketitle
\setcounter{subsection}{1}
\noindent\textit{Editorial note.} All delivered mathematical text is the author's own, copied from the pinned Stacks Project chapter file spaces-morphisms.tex; no statement or proof step was rewritten, reordered or removed.\par\medskip
\noindent\textit{Editorial note.} Uniqueness of the representing morphism is part of the statement; the proof establishes it by showing any two representatives agree on the whole intersection, which is exactly what makes the glued morphism unique.\par\medskip
\noindent\textit{Editorial note.} The single reference of the counted spans is recorded as a non-load-bearing navigation pointer; the argument that consumes it stays inside the counted proof span.\par\medskip
\noindent\textit{Editorial note.} Preamble: the three native excerpt spans of sources/stacks-spaces-morphisms/preamble.tex (lines 8-13, 14-18 and 154-217). The chapter file itself carries no preamble text.\par\medskip
\noindent\textit{Allowed standard prerequisites (editorial).}\begin{itemize}
\item Representatives of a rational map agree on a dense open. used at proof lines 9821-9822. Supplies the dense open W on which the two representatives coincide.
\item The diagonal of a separated morphism is a closed immersion, hence equalizers are closed. used at proof lines 9823-9826. Makes the equalizer E a closed subspace of V intersect V'.
\item A morphism whose target is reduced is determined by its underlying map of topological spaces. Properties of Spaces, Lemma spaces-properties-lemma-map-into-reduction, used at proof lines 9827-9830. Upgrades |E| = |V intersect V'| to an equality of subspaces.
\item Glueing morphisms along a surjection of fppf sheaves. used at proof lines 9831-9834. Produces the global morphism f : U -> Y from the matching local representatives.
\end{itemize}\medskip
\ifdefined\hypertarget\hypertarget{sourcepointer-1}{}\fi
\noindent\textit{Source reference.} Original lemma 03JJ: focused-seeds/AGP/sources/stacks-snapshot-20261011/tags.txt, source lines 4641--4641. Navigation pointer into the Stacks Project; the standard fact it supplies is applied inside the counted proof span and the argument resting on it stays there.\par\medskip
\begin{lemma}
\label{lemma-rational-map-from-reduced-to-separated}
Let $S$ be a scheme. Let $X$ and $Y$ be algebraic spaces over $S$.
Assume $X$ is reduced and $Y$ is separated over $S$. Let
$\varphi$ be a rational map from $X$ to $Y$ with domain of definition
$U \subset X$. Then there exists a unique morphism $f : U \to Y$
of algebraic spaces representing $\varphi$.
\end{lemma}

\begin{proof}
Let $(V, g)$ and $(V', g')$ be representatives of $\varphi$. Then
$g, g'$ agree on a dense open subspace $W \subset V \cap V'$.
On the other hand, the equalizer $E$ of $g|_{V \cap V'}$ and $g'|_{V \cap V'}$
is a closed subspace of $V \cap V'$ because it is the base change
of $\Delta : Y \to Y \times_S Y$ by the morphism
$V \cap V' \to Y \times_S Y$ given by $g|_{V \cap V'}$ and $g'|_{V \cap V'}$.
Now $W \subset E$ implies that $|E| = |V \cap V'|$. As $V \cap V'$
is reduced we conclude $E = V \cap V'$ scheme theoretically, i.e.,
$g|_{V \cap V'} = g'|_{V \cap V'}$, see
Properties of Spaces, Lemma \ref{spaces-properties-lemma-map-into-reduction}.
It follows that we can glue the representatives $g : V \to Y$ of $\varphi$
to a morphism $f : U \to Y$ because $\coprod V \to U$ is a surjection of
fppf sheaves and
$\coprod_{V, V'} V \cap V' = (\coprod V) \times_U (\coprod V)$.
\end{proof}

\end{document}
